\section{Background}
\subsection{API}
Se entende neste texto por API, sigla inglesa para \textit{application
programming interface} como: uma interface que defina nome, entrada e saída de
procedimentos que se referem direta ou indiretamente a uma sequência de
instruções a serem executadas pelo computador.
% preciso de uma citação(?)

É considerado também que tal é normalmente disponibilizado por um grupo com o
intuito de que estes mecanismos sejam utilizados por outros grupos de
desenvolvedores com comunicação relativamente mínima entre todas as partes. O
objetivo principal de tal prática é facilitar o desenvolvimento de software em
contextos / áreas de conhecimento específicas.

Como apontado por\cite{api_property_inference} frequentemente o uso de uma api
pode se resumira uma simples chamada de função, como definido acima, porém
também demanda de seus usuários que uma lógica interna seja seguida para que o
funcionamento seja o esperado pelo usuário da API em questão.

\subsection{Padrões de Uso}
Um padrão de uso na área tem uma definição vaga, variando ligeiramente entre
autores diferentes. Seguindo o proposto em \cite{api_property_inference}, um
padrão de uso é uma forma comum de se utilizar uma API. Utilizando a mesma
fonte, tais padrões podem ser sub categorizados especificamente em padrões de
uso ordenados e não ordenados. No caso dos primeiros, o padrão indica a ordem na
qual as funções são chamadas quando utilizadas em conjunto, enquanto no segundo
caso não há tal restrição.

\subsection{Mineração de Padrões de Uso}

Definição geral do problema Como definidos em textos da
literatura\cite{analyzing_temporal_api_usage_patterns}\cite{api_property_inference},
o processo de mineração de padrões de uso se trata de encontrar em uma base de
dados, de seja códigos fonte, documentações ou mesmo trechos de códigos postados
online em forums, relações entre os métodos, funções e variáveis utilizadas por
desenvolvedores a fim de obter os anteriormente definidos padrões de uso.

\subsection{Grafo de Chamadas}
Um grafo de chamadas trata-se de uma forma de verificar quais funções do código
chamam quais outras funções. Fundamentalmente, se trata de um grafo orientado
onde os vértices são funções, utilizadas ou criadas, um programa. Já as arestas
representam qual função chama qual, no senso de: se A chama B, há uma aresta $A
\rightarrow B$.

Há duas formas principais de obter tal grafo: utilizando traços de execução do
programa ou analisando o código fonte.

O primeiro método envolve inserir uma lógica customizada no código fonte para
registrar em um depósito externo quais funções o programa esta chamando enquanto
está sendo executado. Tal método geralmente demanda um cuidado especial ao
compilar o projeto e, obrigatoriamente, degrada a performance do projeto
original.

Já o segundo método envolve analisar o código fonte com técnicas de análise
sintética. Tal permite ter acesso as funções que cada função chama diretamente
do ponto de vista do programador. Contudo, esta forma de tratar o problema tem a
clara desvantagem de, possivelmente, considerar caminhos de execução que nunca
seriam tomados numa situação real. Adicionalmente, este é incapaz de obter
sequências de chamadas de função indiretas, como poderia ser feito analisando
traços de execução.

Como em qualquer contexto que trata de grafos, há várias possibilidades para
representar o presente tipo grafo. Porém, no caso do grafo de chamadas, o mais
comum é a representação pela lista de adjacência.

\subsection{Minerar Regras de Associação}
Como descrito em \cite{minig_association_rules}, uma regra de associação pode
ser entendida como uma implicação do tipo $X => I_j$ tal que X é um conjunto de
itens em L e $I_j \in L$.

Assim, considera-se que uma transação $t \in T$ satisfaz uma regra de associação
$X => I_j$ se e somente se $t$ possui todos os itens que estão no conjunto $X$ e
o item $I_j$. Dessa forma pode-se medir qual porcentagem $c\%$ do conjunto de
transações $T$ satisfaz uma dada regra de associação, sendo $c$ denominado
fator de confiança.

Um dos primeiros algoritmos a solucionar tal problema de forma satisfatória é
conhecido como apriori\cite{minig_association_rules}. Tal algoritmo consiste em
manter um conjunto de possíveis regras detectadas e, iterativamente, expandir
estes conjuntos em mais um elemento. A fim de evitar minimamente uma explosão
combinatorial, checa-se a cada iteração quais regras estão acima do nível de
confiança aceitável e quais podem ser descartados para a próxima iteração.

\subsection{Clustering}

Como apresentado em
\cite{recommender_system_software_engineering_nearest_neighbours}
\textit{clustering} é uma técnica de aprendizado de máquina não supervisionado
que busca classificar os dados baseado apenas na distribuição destes, sem levar
em conta categorias pré-definidas. Uma vez que tais categorias estão
construídas, novos dados podem ser rapidamente categorizados comparando os com
métricas como o centro dos clusters detectados.

Existem várias técnicas para a obtenção de tais clusters em diferentes tipos de
dados. As técnicas aplicadas dependem principalmente ao tipo de dados no qual
busca-se criar clusters. Comumente, tais técnicas são aplicadas domínio de dados
contínuos, como é o caso de métodos celebres como \textit{K means} e
\textit{Clustering Hierárquico}, disponíveis via interfaces populares e bem
optimizadas como o scipy\cite{scikit_clustering_docs}.

Porém métodos para outros tipos de dados também existem, como o algoritmo
\textit{K modes}\cite{k_modes_article} capaz de encontrar clusters em dados
categóricos.  

Métodos de clustering frequentemente demandam que a quantidade de clusters seja
especificada de antemão. Existem várias heurísticas de como decidir tal valor,
como o método do cotovelo\cite{elbow}, porém em geral não é um problema
solucionável.

\subsection{Modelos Probabilísticos}
Como definido em
\cite{maximum_likelihood_from_incomplete_data_via_the_em_algorithm}, modelos
probabilísticos buscam maximizar a probabilidade do conjunto de dados analisados
segundo seu o modelo que está sendo ajustado. Tal conceito foi expandido através
de processos para ajustar redes Bayesianas\cite{bayesian_network}, onde cada nó
do grafo representa uma variável aleatória e as arestas permitem calcular a
distribuição de probabilidade conjunto das várias variáveis ligadas entre si.

